\documentclass[10pt,a4paper]{amsart}
\usepackage[margin=19mm]{geometry}
\usepackage{amsmath,amssymb,mathtools}
\usepackage[scr=rsfs,cal=euler]{mathalfa}
\usepackage{xfrac}
\usepackage[unicode,hidelinks,bookmarks=false]{hyperref}
\newcommand{\ie}{i.e.\ }
\newcommand{\cf}{cf.\ }
\newcommand{\resp}{respectively\ }
\newcommand{\Subsection}{\subsection}
\providecommand{\nobreakdash}{\nobreak\hbox{-}}
\setlength{\emergencystretch}{2em}
\multlinegap0pt
\usepackage{bm}
\usepackage{bbm}
\usepackage{dsfont}
\usepackage{xspace}
\usepackage{cancel}
\usepackage{mathtools}
\usepackage{amsthm}
\makeatletter
\@ifundefined{theorem}{\newtheorem{theorem}{Theorem}}{}
\@ifundefined{lemma}{\newtheorem{lemma}[theorem]{Lemma}}{}
\makeatother
\makeatletter
\def\FF{{\mathcal F}}
\def\Z{{\mathbb Z}}
\def\bad{{\rm bad}}
\def\cB{{\mathcal B}}
\expandafter\ifx\csname proof\endcsname\relax
\newenvironment{proof}{\noindent {\bf Proof:}}{$\Box$ \vspace{2 ex}}
\fi
\makeatother
\title[Secondary terms in the counting functions of quartic fields]{Secondary terms in the counting functions of quartic fields}
\author{Arul Shankar, Jacob Tsimerman}
\date{Source: arXiv:2503.02381v2 (CC BY 4.0)}
\begin{document}
\maketitle

\noindent\emph{Source and licence.} Secondary terms in the counting functions of quartic fields. Version arXiv:2503.02381v2 (CC BY 4.0), distributed under the Creative Commons Attribution 4.0 International licence, as recorded in the arXiv licence metadata for this exact version and shown on the abstract page https://arxiv.org/abs/2503.02381v2. Full rights notice: licenses/arxiv-2503.02381-CC-BY-4.0.txt.\par\smallskip

\noindent\textit{Editorial scope.} Counted unit: the Lemma whose source label is \texttt{lem:Ab11\_to\_A\_fiber} in the article (source0.tex lines 1252--1259, source label \texttt{lem:Ab11\_to\_A\_fiber}), delivered together with the article's own complete original proof of it (source0.tex lines 1260--1321, one \texttt{proof} environment of 62 lines). No mathematics is written, repaired or re-derived: statement and proof are the article's own text line for line. Required context retained uncounted: thm:projectlargecoords (lines 716--721). Every definition, display and label of those lines is retained unchanged. Omitted from this package and not redistributed: 1--715, 722--1251, 1322--2712. Nothing inside a retained excerpt is omitted or altered: every retained body is delivered exactly as the article's lines are written, so no omission receipt of any kind accompanies this package. Citations: the retained excerpts carry no citation command at all, so no bibliography entry of the article is transcribed and no reference is redistributed. Reference adaptation: none was needed and none was made; every reference inside the delivered unit resolves inside it, so no unresolved reference or citation is produced. Printed slips kept literal: no printed word, symbol, hypothesis or number of the counted statement or of its proof was altered. Layout and class: the article's own class is \texttt{article} with packages including amsfonts, amsmath, amssymb, latexsym, comment, graphicx, xcolor, color, hyperref, titlesec; the wrapper is the lane's pinned \texttt{amsart} editorial wrapper at 10pt on A4, which restates the theorem-like environments the retained text needs and adds the article's own macro definitions that the retained text needs (\texttt{\textbackslash FF}, \texttt{\textbackslash Z}, \texttt{\textbackslash bad}, \texttt{\textbackslash cB}), with extra wrapper packages (bm, bbm, dsfont, xspace, cancel, mathtools). The delivered numbering is the unit's own. Assets: no retained span contains a figure environment, a graphics-inclusion command, a PDF-page inclusion, a table or a TikZ picture; no image file is copied into this package, the package declares no asset dependency and no image file is redistributed with it. Licence: version arXiv:2503.02381v2 of this article is distributed under the Creative Commons Attribution 4.0 International licence, as recorded in the arXiv licence metadata for this exact version and shown on the abstract page https://arxiv.org/abs/2503.02381v2. A full rights notice is delivered at licenses/arxiv-2503.02381-CC-BY-4.0.txt. Characters: every retained excerpt is pure ASCII, so the delivered engine, XeLaTeX with its default Latin Modern text font, reports no missing character.
\begin{theorem}\label{thm:projectlargecoords}
    Let the notation be as above, with $f,B$ fixed and all other parameters varying. Suppose for some parameter $Y$ we have $d_1,\dots,d_r\geq Y$ and $d_{s+1},\dots,d_n\leq Y^{-1}$. Then
\begin{equation*}
\sum_{\ell\in \Z^n}(gB)(\ell)f(\ell) = \prod_{i=1}^r d_i  \sum_{\ell_0\in \Z^{s-r}}(g_{r,s}B_{r,s})(\ell_0)f_{r,s}(\ell_0) + O_A(Y^{-A}).
\end{equation*}
\end{theorem}

\begin{lemma}\label{lem:Ab11_to_A_fiber}
We have
\begin{equation*}
\int_{g\in\FF}\bigl(E_g^{(1)}-E_g^{(2)}\bigr)
\psi\Big(\frac{\lambda(g)}{X^{1/12}}\Big)f^X(g)dg
=O(X^{4/5+o(1)}).
\end{equation*}
\end{lemma}

\begin{proof}
  Let $\theta<\delta$ be a small constant to be picked later. Note that if $g\in\FF_3$ with $w_g(b_{11})\gg X^\theta$, then both $E_g^{(1)}$ and $E_g^{(2)}$ are superpolynomially small by Theorem \ref{thm:projectlargecoords}. Hence, we may restrict the integral above to $g\in\FF^\bad$, where $\FF^\bad$ consists of $g\in\FF$ with $f^X(g)>0$ and $w_g(b_{11})\ll X^\theta$.

%Furthermore, if $w_g(b_{11})$, $w_g(a_{22})$, or $w_g(a_{13})$ are less than $c$ for some sufficiently small constant $c$, then the sum over $L$ is empty (since $L$ contains no elements with $a_{11}=b_{11}=0$ or $\det(A)=0$). 
%Thus, we may restrict the integral above to $\FF^{\bad}$, the region of $g\in\FF$ for which $f^X(g)>0$, $\lambda(g)\asymp X^{1/12}$, $1\ll w_g(a_{13})$, $1\ll w_g(a_{22})$, and $1\ll w_g(b_{11})\ll X^\theta$.

Next, we note that $E_g^{(1)}-E_g^{(2)}=F_g^{(1)}-F_g^{(2)}$, where 
\begin{equation*}
\begin{array}{rcl}
F^{(1)}_g&:=&\displaystyle\sum_{A,b_{11}}\nu(L|_{A,b_{11}})(g\cB)_{\emptyset;\{A,b_{11}\}}(A,b_{11})-\sum_{b_{11}}\nu(L|_{a_{11}=0;b_{11}})(g\cB)_{\{a_{11}\};\{b_{11}\}}(b_{11});
\\[.15in]
F^{(2)}_g&:=&\displaystyle\sum_{A}\nu(L_{A})(g\cB)_{\emptyset;\{A\}}(A)-\nu(L_{a_{11}=0})(g\cB)_{\{a_{11}\}}.
\end{array}
\end{equation*}
Define $\FF^{\bad,1}$ to be the set of $g\in\FF^{\bad}$ satisfying $w_g(a_{12})\ll X^\theta$. We claim that for $g\in\FF^\bad\backslash\FF^{\bad,1}$, we have $F_g^{(1)}-F_g^{(2)}$ is superpolynomially small: 
Indeed, when $w_g(a_{12})>X^\theta$, we have upto superpolynomially small error,
\begin{equation*}
F_g^{(1)}-F_g^{(2)}=\sum_{a_{11}\neq 0,b_{11}}
\nu(L|_{a_{11},b_{11}})(g\cB)_{\emptyset;\{a_{11},b_{11}\}}
-\sum_{a_{11}\neq 0}
\nu(L|_{a_{11}})(g\cB)_{\emptyset;\{a_{11}\}}.
\end{equation*}
Since $a_{11}\neq 0$ and $t\gg X^\delta$ forces the range of $b_{11}$ to be large, the above is superpolynomially small.
Thus, up to superpolynomially small error, we may restrict the integral over $g\in\FF$ in the displayed equation of the lemma to $g\in\FF^{\bad,1}$. 


We now show that the integral over $g\in\FF^{\bad,1}$ of all four
terms constituting $E_g^{(1)}$ and $E_g^{(2)}$ are small.
Let us begin
with the first term of $E_g^{(1)}$: Note that if any of $w_g(b_{11})$, $w_g(a_{12})$, or $w_g(a_{22})$ are less than a sufficiently small positive constant $c$, then
$\nu(L_{A,b_{11}}))(g\cB_{\phi;\{A,b_{11}\}})$ is $0$. This is because $L$ contains no elements with $\det(A)=0$ or $(a_{11},b_{11})=(0,0)$.
Note also that since $g\in\FF^{\bad,1}$ is such that $w_g(b_{11})<X^\theta$ and $f^X(g)>0$, it follows that $g\cB$ is nonzero only on integral pairs $(A,B)$ with $a_{11}=0$. Hence we have
\begin{equation}\label{eq:badbound1}
\begin{array}{rcl}
  \displaystyle  \int_{\substack{g\in\FF^{\bad,1}}}\sum_{A,b_{11}}\nu(L|_{A,b_{11}})(g\cB)_{\emptyset;\{A,b_{11}\}}\psi\Bigl(
\frac{\lambda(g)}{X^{1/12}}
\Bigr)dg&\ll& \displaystyle\int_{\substack{\FF^{\bad,1}\\\lambda(g)\asymp X^{1/12}}}\max(1,w_g(a_{12}))\prod_{\alpha\not\in\{a_{11},a_{12}\}}w_g(\alpha)
\\[.2in]&\ll&\displaystyle
X^{5/6+\theta}\int_{\substack{(t,s_1,s_2)\in\FF^{\bad,1}}}s_1^{-1}s_2^{-2}d^\times td^\times s_1 d^\times s_2.
\end{array}
\end{equation}
In fact, the same bound is true for each of the four terms
constituting $E_g^{(1)}$ and $E_g^{(2)}$. For example, for the second
term of $E_g^{(1)}$, though we can no longer assume that
$w_g(b_{11})>c$, the estimate remains the same since only the volume of the $b_{11}$-projection is being integrated (and not the number of integral choices for $b_{11}$).

%The conditions $w_g(a_{13})\gg 1$, $w_g(a_{22})\gg 1$, and $w_g(b_{11})\gg 1$ imply that $t$, $s_1$, and $s_2$ are bounded by some fixed power of $\lambda$ for all $\lambda(t,s_1,s_2)\in \FF^\bad$. Moreover,
The conditions $w_g(b_{11})\ll X^\theta$ and $w_g(a_{12})\ll X^\theta$
respectively imply that we have
\begin{equation*}
Y:=\frac{s_1^4s_2^2 X^\theta}{t\lambda}\gg 1;\quad
Y':=\frac{ts_1s_2^2 X^\theta}{\lambda}\gg 1.
\end{equation*}
Consider the integral in the final line of
\eqref{eq:badbound1}. Multiplying the integrand
by $Y^{1/5+\epsilon}Y'^{1/5-5\epsilon}$, yields
\begin{equation*}
  X^{5/6+\theta}\int_{\substack{(t,s_1,s_2)\in\FF^{\bad,1}}}s_1^{-1}s_2^{-2}d^\times td^\times s_1 d^\times s_2
\ll X^{4/5+O(\theta+o(1))}.
\end{equation*}
Since $\theta$ and $\epsilon$ can be taken to be arbitrarily small, the result follows.
\end{proof}

\end{document}
